% Created 2024-01-08 lun 09:33
% Intended LaTeX compiler: pdflatex
\documentclass[11pt]{article}
\usepackage[utf8]{inputenc}
\usepackage[T1]{fontenc}
\usepackage{graphicx}
\usepackage{longtable}
\usepackage{wrapfig}
\usepackage{rotating}
\usepackage[normalem]{ulem}
\usepackage{amsmath}
\usepackage{amssymb}
\usepackage{capt-of}
\usepackage{hyperref}
\usepackage{minted}

\usepackage{parskip}
\usepackage[utf8]{inputenc} %% For unicode chars
\usepackage{placeins}
\usepackage[margin=2.50cm]{geometry}
\usepackage[T1]{fontenc}
\usepackage{mathpazo}
\linespread{1.05}
\usepackage[scaled]{helvet}
\usepackage{courier}
\hypersetup{colorlinks=true,linkcolor=blue}
\RequirePackage{fancyvrb}
\DefineVerbatimEnvironment{verbatim}{Verbatim}{fontsize=\small,formatcom = {\color[rgb]{0.5,0,0}}}
\usepackage{mdframed}
\BeforeBeginEnvironment{minted}{\begin{mdframed}}
\AfterEndEnvironment{minted}{\end{mdframed}}
\date{}
\title{}
\hypersetup{
 pdfauthor={},
 pdftitle={},
 pdfkeywords={},
 pdfsubject={},
 pdfcreator={Emacs 28.1 (Org mode 9.5.5)}, 
 pdflang={Spanish}}
\begin{document}

\setlength\parindent{10pt}

\section*{User management}
\label{sec:org90866e5}
Creating a \texttt{User} class in PHP involves implementing methods for user
authentication, role-based authorization, and session management. Below
is an outline of the \texttt{User} class, including methods for login,
registration, role checking, and session handling. Additionally, I'll
provide examples of corresponding HTML forms for user interaction.

\subsection*{PHP \texttt{User} Class}
\label{sec:orgd6c866b}
\begin{minted}[]{php}
class User {
    private $pdo;
    private $userId;
    private $username;
    private $role;

    public function __construct($pdo) {
        $this->pdo = $pdo;
        session_start();
    }

    public function login($username, $password) {
        $stmt = $this->pdo->prepare("SELECT id, username, password, role FROM users WHERE username = ?");
        $stmt->execute([$username]);
        $user = $stmt->fetch();

        if ($user && password_verify($password, $user['password'])) {
            $_SESSION['user_id'] = $user['id'];
            $_SESSION['username'] = $user['username'];
            $_SESSION['role'] = $user['role'];
            $this->userId = $user['id'];
            $this->username = $user['username'];
            $this->role = $user['role'];
            return true;
        }
        return false;
    }

    public function register($username, $password) {
        $hashedPassword = password_hash($password, PASSWORD_DEFAULT);
        $stmt = $this->pdo->prepare("INSERT INTO users (username, password) VALUES (?, ?)");
        $stmt->execute([$username, $hashedPassword]);
    }

    public function logout() {
        session_unset();
        session_destroy();
    }

    public function isUserLoggedIn() {
        return isset($_SESSION['user_id']);
    }

    public function getUsername() {
        return $this->username;
    }

    public function getRole() {
        return $this->role;
    }

    public function isAdmin() {
        return $this->role === 'admin';
    }
}
\end{minted}

\subsection*{Explanation}
\label{sec:orgc6d32ed}
\begin{itemize}
\item \textbf{Constructor}: Initializes the database connection and starts the
session.
\item \textbf{Login}: Verifies user credentials and starts a session for the user.
\item \textbf{Register}: Registers a new user with a hashed password.
\item \textbf{Logout}: Ends the session.
\item \textbf{Session Checks}: Methods like \texttt{isUserLoggedIn} and \texttt{isAdmin} check
the user's login status and role.
\end{itemize}

\subsection*{HTML Forms}
\label{sec:org5b5f2bf}
\subsubsection*{Registration Form (\texttt{register.html})}
\label{sec:org275d539}
\begin{minted}[]{php}
<!DOCTYPE html>
<html>
<head>
    <title>Register</title>
</head>
<body>
    <h1>Register</h1>
    <form action="register.php" method="post">
        Username: <input type="text" name="username" required><br>
        Password: <input type="password" name="password" required><br>
        <input type="submit" value="Register">
    </form>
</body>
</html>
\end{minted}

\subsubsection*{Login Form (\texttt{login.html})}
\label{sec:org83b4fb1}
\begin{minted}[]{html}
<!DOCTYPE html>
<html>
<head>
    <title>Login</title>
</head>
<body>
    <h1>Login</h1>
    <form action="login.php" method="post">
        Username: <input type="text" name="username" required><br>
        Password: <input type="password" name="password" required><br>
        <input type="submit" value="Login">
    </form>
</body>
</html>
\end{minted}

\subsection*{PHP Scripts for Forms}
\label{sec:org9ef2a00}
\subsubsection*{Registration Script (\texttt{register.php})}
\label{sec:org8a84ccd}
\begin{minted}[]{php}
<?php
require_once 'User.php';
require_once 'db_connect.php'; // PDO connection

$user = new User($pdo);

if ($_SERVER["REQUEST_METHOD"] == "POST") {
    $user->register($_POST['username'], $_POST['password']);
    // Redirect to login or display a success message
}
\end{minted}

\subsubsection*{Login Script (\texttt{login.php})}
\label{sec:orgf70c0b3}
\begin{minted}[]{php}
<?php
require_once 'User.php';
require_once 'db_connect.php'; // PDO connection

$user = new User($pdo);

if ($_SERVER["REQUEST_METHOD"] == "POST") {
    if ($user->login($_POST['username'], $_POST['password'])) {
        header("Location: dashboard.php"); // Redirect to a protected page
    } else {
        echo "Invalid credentials.";
    }
}
\end{minted}

\subsection*{Security Considerations}
\label{sec:org123a5a1}
\begin{itemize}
\item Always sanitize and validate user inputs.
\item Store hashed passwords in the database, never plain text.
\item Implement CSRF protection on forms.
\item Ensure secure session management, particularly on shared hosting
environments.
\end{itemize}

This setup provides a foundational structure for user authentication and
role-based authorization in your PHP application, integrating session
management for maintaining user state across different pages.
\end{document}
